\chapter{Academic Rubric Traceability}
\label{app:rubric}

\section{Compliance matrix}

This appendix maps the 2025--2026 report instructions and grading rubric to the
corresponding sections of the document.

\begin{longtable}{@{}p{0.30\textwidth}p{0.56\textwidth}@{}}
\caption{Traceability to the report instructions and rubric.}
\label{tab:rubric-trace}\\
\toprule
\textbf{Criterion} & \textbf{Location/evidence}\\
\midrule
\endfirsthead
\toprule
\textbf{Criterion} & \textbf{Location/evidence}\\
\midrule
\endhead
Cover and identification & Cover page and identification sheet include
student, programme, host, supervisors and internship dates.\\
Summary and keywords & English abstract and French résumé precede the table of
contents; both include keywords.\\
Acknowledgements & Dedicated unnumbered section in the front matter.\\
Glossary & Acronym list defines ROS, Gazebo/PDDL/PlanSys2, QoS and ML, and the
external VLA architecture discussed in the related work
terms used in the body.\\
Table of contents and lists & Detailed contents plus lists of figures and
tables.\\
Numbering and pagination & Chapters, sections, figures, tables and references
are numbered; Arabic-numeral pagination begins at the cover, whose visible
number is suppressed.\\
Introduction & \Cref{chap:introduction} states context, engineering objective,
method, contributions, evidence boundary and structure.\\
Host organisation & \Cref{chap:host} covers legal/sector data, APE 85.42Z,
size, organisation, MFJA service, tutors and the student's position within the
six-page limit.\\
Mission and project management & \Cref{chap:mission} provides the administrative
title, the initial Room~315 assignment, scope categories, work packages,
schedule, tools, risks, configuration, communication and change history.\\
Technical context and method & \Cref{chap:background}--\cref{chap:interfaces}
explain ROS~2/Gazebo, digital twin, robots, rail, typed interfaces and safety.\\
Complete engineering scope & \Cref{chap:digital-twin}--\cref{chap:interfaces}
cover the floor, four robot families, unified interface and transport before
the learned work.\\
Learned and deterministic roles & \Cref{chap:ai-development} explains the
current visual-fact boundary and the allocation of validation, planning,
supervision, controller and feedback authority.\\
Current system & \Cref{chap:neurosymbolic} describes language, observation,
PlanSys2, supervision, effect verification and topology-aware recovery.\\
Results and interpretation & \Cref{chap:ai-development,chap:results} report
the current visual-data contract, configured-checkpoint validation, layered
automated evidence and the successful 12-case manifest-bound integrated
campaign.\\
Perspectives & \Cref{chap:perspectives} separates current technical limits,
prioritised validation, continued MFJA use and focused simulation research
extensions.\\
Engineering qualities & \Cref{chap:mission} identifies personal work,
initiative, autonomy, supervisory collaboration, resource management and
professional learning through two reflective engineering examples and a
competency synthesis.\\
Conclusion & \Cref{chap:conclusion} synthesises objectives, results,
limitations, learning and next steps without introducing a new result.\\
Illustrations & Figures are captioned, numbered, referenced and explained in
the text; project captures and explanatory diagrams are identified.\\
Bibliography/webography & Primary standards, official host data, ROS/Gazebo
documentation and research papers are cited with BibTeX.\\
Appendices & Reproducibility, interface, verification and rubric-traceability
material support the main analysis.\\
Report language & English.\\
\bottomrule
\end{longtable}

\section{Length and synthesis}

The instructions require at least 40 substantive pages, excluding appendices
and source listings. The main body covers the project context, architecture,
implementation, experimental evidence and limitations. Generated code and raw
logs remain outside the main discussion, while tables and figures are used for
comparisons, system responsibilities and execution sequences.
